\documentclass[aspectratio=169,10pt,t]{beamer}
\usepackage{issa}  % place issa.sty in the same folder
\definecolor{AccentCol}{HTML}{A8861A}
\renewcommand{\lecnum}{10}
\begin{document}
\titleslide{10}{Project Presentations \&\\Q\&A Panel}{Applying AI \& ML to a defence problem of your choosing\\Structure \textbullet\ rubric \textbullet\ project catalogue \textbullet\ Q\&A panel guide \textbullet\ course conclusion}

\sect{1. Session}
\begin{frame}{Lecture 10 --- How the Session Runs}
\figc{1.6cm}{\begin{tikzpicture}[x=0.075cm]
\fill[AccentCol!80](0,0)rectangle(60,0.8);\fill[ISSAnavy!80](60,0)rectangle(120,0.8);\fill[ISSAgrn!75](120,0)rectangle(180,0.8);
\node[white,font=\scriptsize\bfseries]at(30,0.4){HOUR 1: presentations};\node[white,font=\scriptsize\bfseries]at(90,0.4){HOUR 2: panel deep-dive Q\&A};\node[white,font=\scriptsize\bfseries]at(150,0.4){HOUR 3: course conclusion};
\foreach \x in {0,12,24,36,48}{\draw[white,thick](\x,0)--(\x,0.8);}\node[lab]at(30,-0.35){5 groups $\times$ (10 min talk + 2 min Q\&A)};
\end{tikzpicture}}
\begin{columns}[T]
\column{0.50\textwidth}
\begin{cbox}{Schedule (3 hours)}\scriptsize\textbf{Hour 1:} each group of 4--5 students gives a 10-minute presentation + 2 minutes of panel Q\&A; slots drawn by lot; 4--6 groups expected. \textbf{Hour 2:} open floor --- any question on any lecture. Panel: course instructor + 1--2 external defence-AI practitioners; focus on how your algorithm would perform in REAL operational conditions. \textbf{Hour 3:} capstone feedback summary, industry perspectives, recommended reading.\end{cbox}
\column{0.48\textwidth}
\begin{cbox}[AccentCol]{What the panel is looking for}\scriptsize Think like a defence-AI practitioner, not a textbook. \textbf{\textcircled{\tiny1} Problem framing:} the right problem type, a justified algorithm choice. \textbf{\textcircled{\tiny2} Honest limitations:} sim-to-real gap quantified; data pitfalls acknowledged. \textbf{\textcircled{\tiny3} Operational applicability:} deployment, autonomy level, ROE. \textbf{\textcircled{\tiny4} Critical evaluation:} would a sceptical senior officer accept your 94\% simulation result?\end{cbox}
\end{columns}
\end{frame}

\sect{2. Presentation}
\begin{frame}{The Ten-Minute Framework --- Six Sections, One Story}
\figc{1.3cm}{\begin{tikzpicture}[x=0.024cm]
\foreach \a/\b/\t/\cl in {0/60/1 Problem/ISSAnavy,60/120/2 Data/ISSAblue,120/210/3 Algorithm/AccentCol,210/330/4 Results/ISSAgrn,330/420/5 Deployment/violet,420/450/6 Limits/ISSAred}{\fill[\cl!75](\a,0)rectangle(\b,0.6);\node[white,font=\fontsize{5.5}{6}\selectfont\bfseries]at({(\a+\b)/2},0.3){\t};}
\foreach \x/\l in {0/0,60/1:00,120/2:00,210/3:30,330/5:30,420/7:00,450/7:30}{\node[font=\tiny]at(\x,-0.2){\l};}
\end{tikzpicture}}
\sm{The six timed sections fill 7.5 minutes; the rest of the 10-minute slot is margin for transitions and a demonstration.}
\figc{4.0cm}{\begin{tikzpicture}
\tikzset{P/.style={nd,text width=4.75cm,align=left,minimum height=1.75cm,font=\fontsize{6.3}{7.5}\selectfont,inner sep=3pt}}
\node[P,draw=ISSAnavy]at(-5,0.95){\textbf{1. Problem (60 s).} Not ``we applied ML to defence'' but: ``IADS operators receive 400 tracks/hour; manual classification takes 45 s each; we need $\ge$95\% accuracy in $\le$3 ms.'' Name the threat environment; quantify the cost of NOT solving it.};
\node[P,draw=ISSAblue]at(0,0.95){\textbf{2. Data \& quality (60 s).} $N$, features, class distribution --- and which L1 pitfalls you found and fixed (shift, label noise, imbalance, confounders). Never mentioning data quality is a red flag.};
\node[P,draw=AccentCol]at(5,0.95){\textbf{3. Algorithm (90 s).} Which and WHY, against at least one alternative: ``RF over SVM because 8,500 samples $\times$ 11 features favour ensemble variance reduction over kernel sensitivity.'' Cite course benchmarks.};
\node[P,draw=ISSAgrn]at(-5,-0.95){\textbf{4. Results (120 s).} Headline metric, then the breakdown: per-class recall, confusion matrix, silhouette, AUROC. Address imbalance; one chart comparing your result with a baseline.};
\node[P,draw=violet]at(0,-0.95){\textbf{5. Deployment (90 s).} Autonomy level, latency, human--AI interface, ROE compliance, estimated sim-to-real gap and how you would validate before fielding.};
\node[P,draw=ISSAred]at(5,-0.95){\textbf{6. Limitations (30 s).} TWO specific, honest ones: ``trained only on L-band signatures; frequency-agile X-band radars need retraining.'' Panels respect honesty over false confidence.};
\end{tikzpicture}}
\end{frame}

\sect{3. Rubric}
\begin{frame}{Assessment Rubric --- Five Criteria, Each Marked Out of 20}
\figc{3.6cm}{\scriptsize\renewcommand{\arraystretch}{1.15}\begin{tabular}{p{2.4cm}cp{2.9cm}p{2.3cm}p{2.1cm}p{1.8cm}}\toprule\textbf{Criterion}&\textbf{Weight}&\textbf{Distinction (18--20)}&\textbf{Merit (14--17)}&\textbf{Pass (10--13)}&\textbf{Fail ($<$10)}\\\midrule Problem definition \& context&20\%&precisely quantified; operationally grounded; cost of inaction&clear, most context&lacks operational grounding&vague, generic\\Algorithm justification&25\%&compared with 2+ alternatives; cites course metrics&one alternative&no clear justification&wrong problem type\\Results \& validation&25\%&per-class metrics, CIs, method, baseline&several metrics&primary metric only&no context or baseline\\Operational applicability&20\%&autonomy, latency, ROE, sim-to-real gap, plan&most factors&gaps remain&purely academic\\Critical evaluation \& ethics&10\%&limitations, IHL, adversarial robustness, failure modes&two limitations&one limitation&none\\\bottomrule\end{tabular}}
\begin{columns}[T]
\column{0.50\textwidth}
\figc{1.9cm}{\bars{25}{1.3}{1.6}{Problem/20/ISSAnavy,Algorithm/25/AccentCol,Results/25/ISSAgrn,Operational/20/violet,Ethics/10/ISSAred}}
\sm{Criterion weights (\%).}
\column{0.48\textwidth}
\begin{ibox}\scriptsize\textbf{Peer evaluation} contributes 20\% of the project grade: each student rates OTHER groups on a 5-point scale for clarity, technical depth, operational relevance and question handling; moderated by the instructor.\end{ibox}
\end{columns}
\end{frame}

\sect{4. Topics}
\begin{frame}{Project Topic Catalogue --- Eight Suggested Themes}
\figc{5.2cm}{\begin{tikzpicture}
\tikzset{T/.style={nd,text width=3.45cm,align=left,minimum height=2.45cm,font=\fontsize{6.1}{7.2}\selectfont,inner sep=3pt}}
\node[T,draw=ISSAgrn,fill=ISSAgrn!6]at(-5.4,1.3){\textbf{T1 Threat classification (L2).} IADS classifier for a new sensor (EO/IR, SAR, SIGINT); compare RF, XGBoost, SVM; fix imbalance; pick a platform. \textbf{Benchmark:} RF 97.5\%.};
\node[T,draw=violet,fill=violet!6]at(-1.8,1.3){\textbf{T2 Mission-planning regression (L3).} Predict fuel, duration, munitions or rescue time from JANUS data; OLS vs Ridge vs XGBoost; fuse $\ge$2 inputs. \textbf{Benchmark:} $R^2$ 0.912.};
\node[T,draw=teal,fill=teal!6]at(1.8,1.3){\textbf{T3 OPFOR pattern intelligence (L4).} Cluster movement/comms/logistics data with $k$-Means + DBSCAN + PCA; find one unexpected cluster. \textbf{Benchmark:} silhouette 0.712.};
\node[T,draw=ISSAred,fill=ISSAred!5]at(5.4,1.3){\textbf{T4 Cyber-defence anomaly (L5).} Tiered IF + AE for a military network; test on one novel attack; compute daily load. \textbf{Benchmark:} AE TPR 95.8\%.};
\node[T,draw=blue!60!black,fill=blue!5]at(-5.4,-1.3){\textbf{T5 Wargame decision support (L6).} DQN advisor with shaping and curriculum; win rate vs random and vs human; Q-value output. \textbf{Benchmark:} 93\% vs random.};
\node[T,draw=teal!60!black,fill=green!6]at(-1.8,-1.3){\textbf{T6 Autonomous platform control (L7).} PPO or SAC for a drone, interceptor, ship autopilot or submarine evasion; domain randomisation. \textbf{Benchmark:} SAC 96.2\%.};
\node[T,draw=orange!70!black,fill=orange!6]at(1.8,-1.3){\textbf{T7 Multi-agent operations (L8).} MARL for 3+ agents with CTDE (QMIX or MAPPO); identify an emergent behaviour. \textbf{Benchmark:} CommsMAPPO 88.1\%.};
\node[T,draw=AccentCol,fill=AccentCol!8]at(5.4,-1.3){\textbf{T8 Original proposal.} Combine 2+ lectures; approved by Session 8; novel data or scenario, clear operational context, one course baseline.};
\end{tikzpicture}}
\end{frame}

\sect{5. Q\&A Panel}
\begin{frame}{The Q\&A Panel --- It Tests Judgement, Not Algorithm Recall}
\begin{columns}[T]
\column{0.50\textwidth}
\begin{cbox}[AccentCol]{Every question probes one of three things}\scriptsize \textbf{\textcircled{\tiny1} Depth:} ``You reported 94\% accuracy. What is the AUROC? The confusion matrix for the rarest class?'' \textbf{\textcircled{\tiny2} Judgement:} ``Your DQN won 87\% in simulation. Command is considering Level-3 deployment. Your recommendation?'' \textbf{\textcircled{\tiny3} Intellectual honesty:} ``Give me a specific scenario where your model would fail.''\end{cbox}
\figc{2.0cm}{\begin{tikzpicture}\node[ndA,minimum width=2.0cm](d)at(0,0){DEPTH};\node[ndN,minimum width=2.0cm](j)at(2.4,0){JUDGEMENT};\node[ndR,minimum width=2.0cm](h)at(4.8,0){HONESTY};\node[lab]at(2.4,-0.7){beyond the headline \quad real-world implications \quad failure modes};\end{tikzpicture}}
\column{0.48\textwidth}
\begin{cbox}[ISSAnavy]{Anticipated questions by topic}\scriptsize \textbf{Data (L1):} ``How did you rule out survivorship bias?'' ``What if the enemy changes behaviour to fool your classifier?'' \textbf{Performance (L2--3):} ``What is the operational cost of a false positive?'' ``Show me per-class recall.'' \textbf{Clustering/anomaly (L4--5):} ``How do you know the clusters are real? ARI?'' ``2\% FPR on 1M packets/day is 20,000 false alerts --- how does an analyst cope?'' \textbf{RL/MARL (L6--8):} ``What is the sim-to-real gap?'' ``Who is accountable if the agent errs?''\end{cbox}
\end{columns}
\end{frame}

\begin{frame}{Handling Tough Questions --- the STARE Framework}
\figc{2.3cm}{\begin{tikzpicture}
\tikzset{L/.style={circle,draw=AccentCol,fill=AccentCol!20,minimum size=0.8cm,font=\bfseries},D/.style={font=\fontsize{6}{7}\selectfont,text width=2.3cm,align=center,anchor=north}}
\foreach \i/\l/\t in {0/S/STATE what you DO know,1/T/TIE it to your project slides and results,2/A/ACKNOWLEDGE the gap honestly --- never bluff,3/R/REASON about what would happen,4/E/END with what you would do next}{\node[L]at(\i*2.6,0){\l};\node[D]at(\i*2.6,-0.5){\t};}
\foreach \i in {0,1,2,3}{\draw[arr](\i*2.6+0.42,0)--(\i*2.6+2.18,0);}
\end{tikzpicture}}
\begin{columns}[T]
\column{0.50\textwidth}
\begin{cbox}[ISSAgrn]{Worked STARE answer}\scriptsize\textit{Q: ``Is your detector robust to adversarial inputs?''}\\ \textbf{S:} ``Validated on 100,000 packets: TPR 95\%, FPR 0.8\%.'' \textbf{T:} ``Slide 4 shows the per-attack breakdown.'' \textbf{A:} ``We did not test adversarial perturbations.'' \textbf{R:} ``Autoencoders can be evaded by attacks that mimic normal reconstruction, as R2L partly does.'' \textbf{E:} ``Next: red-team with crafted traffic and add an Isolation-Forest cross-check.''\end{cbox}
\column{0.48\textwidth}
\begin{cbox}[AccentCol]{Five phrases that impress panels}\scriptsize\begin{te}\item ``We benchmarked against Lecture $N$'s [X\%] and achieved [Y\%]: a [Z]-point difference due to [reason].''\item ``We chose [A] over [B] because our data had [property] that favours [A].''\item ``We estimate a sim-to-real gap of $\sim$[X]\% from L7 (PPO 94\% $\to$ 78\%) and designed domain randomisation to close it.''\item ``Our rarest class has recall [R]; we set the threshold for the cost of a miss.''\item ``We recommend Level 2: the system advises, the commander decides.''\end{te}\end{cbox}
\end{columns}
\end{frame}

\sect{6. Pitfalls}
\begin{frame}{Common Project Pitfalls --- Avoid These in Your Presentation}
\figc{4.6cm}{\begin{tikzpicture}
\tikzset{M/.style={nd,text width=3.45cm,align=left,minimum height=2.05cm,font=\fontsize{6.2}{7.3}\selectfont,inner sep=3pt,draw=ISSAred,fill=ISSAred!4}}
\node[M]at(-5.4,1.15){\textbf{1 Accuracy-only reporting.} 94\% without confusion matrix, per-class recall or AUC is misleading on imbalanced data.};
\node[M]at(-1.8,1.15){\textbf{2 No baseline.} Compare with the majority class, the course benchmark and an alternative algorithm --- 87\% means nothing alone.};
\node[M]at(1.8,1.15){\textbf{3 Ignoring data quality.} Every real dataset has bias, shift or label noise; skipping L1 is a red flag.};
\node[M]at(5.4,1.15){\textbf{4 Simulation = readiness.} ``Ready to deploy'' from simulation alone; the sim-to-real gap is never zero.};
\node[M]at(-5.4,-1.15){\textbf{5 Missing ethics.} No autonomy level, ROE, IHL, accountability chain or failure consequences.};
\node[M]at(-1.8,-1.15){\textbf{6 Wrong algorithm.} Deep learning by default when a simpler model works; supervised methods with no labels.};
\node[M]at(1.8,-1.15){\textbf{7 Misread statistics.} Adding percentages, ignoring base rates, ``not significant = no effect'', correlation as cause (L1 Section 0).};
\node[M]at(5.4,-1.15){\textbf{8 No uncertainty.} Point results without confidence intervals or repeated runs; a 1-point lead may be noise.};
\end{tikzpicture}}
\sm{Pitfalls 1--6 are from the original lecture pack; 7--8 add the Lecture 1 data-interpretation lessons.}
\end{frame}

\sect{7. Across Disciplines}
\begin{frame}{Presenting Evidence Responsibly --- Lessons From Other Disciplines}
\figc{5.0cm}{\begin{tikzpicture}
\tikzset{G/.style={nd,text width=4.75cm,align=left,minimum height=2.3cm,font=\fontsize{6.9}{8.2}\selectfont,inner sep=3pt}}
\node[G,draw=ISSAgrn,fill=ISSAgrn!7]at(-5,1.25){\textbf{\textcolor{ISSAgrn}{MEDICINE}}\\Clinical trials report to the CONSORT standard: how many enrolled, how many dropped out, the pre-specified outcome, effect size with confidence interval.\\\textbf{Your project:} report $N$, exclusions, the primary metric fixed in advance, and CIs.};
\node[G,draw=ISSAblue,fill=ISSAblue!7]at(0,1.25){\textbf{\textcolor{ISSAblue}{AERONAUTICS}}\\Aviation's checklist and crew-resource-management culture lets any crew member challenge the captain --- decisions are verified, not assumed.\\\textbf{Your project:} a pre-deployment checklist and a clear human override.};
\node[G,draw=ISSAred,fill=ISSAred!6]at(5,1.25){\textbf{\textcolor{ISSAred}{ENGINEERING}}\\Design reviews use failure-mode and effects analysis (FMEA): list how the system can fail, how likely, how severe, and how it is detected.\\\textbf{Your project:} a failure-mode table for your model.};
\node[G,draw=violet,fill=violet!7]at(-5,-1.25){\textbf{\textcolor{violet}{PSYCHOLOGY}}\\After the replication crisis, pre-registration --- publishing hypotheses and analysis plans before seeing the data --- became good practice.\\\textbf{Your project:} state your hypothesis and metric before tuning.};
\node[G,draw=ISSAamb,fill=ISSAyell!45]at(0,-1.25){\textbf{\textcolor{ISSAamb!80!black}{STATISTICS}}\\Reproducibility: fixed random seeds, shared code and data, repeated runs with variance reported.\\\textbf{Your project:} report mean $\pm$ standard deviation over at least 5 seeds for RL results.};
\node[G,draw=AccentCol,fill=AccentCol!8]at(5,-1.25){\textbf{\textcolor{AccentCol}{MILITARY PRACTICE}}\\Red teaming and after-action reviews: an independent team attacks the plan; lessons are recorded without blame.\\\textbf{Your project:} let another group red-team your model before the panel does.};
\end{tikzpicture}}
\end{frame}

\sect{8. Conclusion}
\begin{frame}{Course Conclusion --- The Arc of Ten Lectures}
\begin{columns}[T]
\column{0.52\textwidth}
\figc{3.0cm}{\begin{tikzpicture}
\foreach \i/\t/\cl in {0/{Data can lie: L1}/AccentCol,1/{Learn from labels: L2--L3}/ISSAgrn,2/{Discover without labels: L4--L5}/teal,3/{Decide by interaction: L6--L8}/blue!60!black,4/{Integrate \& present: L9--L10}/violet}{\node[anchor=west,draw=\cl,fill=\cl!12,minimum height=0.5cm,minimum width={2.8cm+\i*0.7cm},font=\fontsize{6.4}{7}\selectfont\bfseries,inner sep=2pt]at(0,\i*0.6){\t};}
\end{tikzpicture}}
\begin{cbox}[ISSAnavy]{The habit that matters most}\scriptsize Ask of every dataset and every AI output: \emph{Who collected this? Under what conditions? What is missing? What changed between collection and deployment? How would I know if it were wrong?} That question will save lives.\end{cbox}
\column{0.46\textwidth}
\begin{rbox}\scriptsize\textbf{What you can now do:} frame defence problems as ML tasks; choose and justify algorithms with numbers; validate honestly; interpret results without the classic statistical errors; design human--machine teaming that keeps humans accountable.\end{rbox}
\begin{wbox}\scriptsize AI will not replace military commanders --- but commanders who understand AI's strengths \emph{and its failure modes} will lead those who do not.\end{wbox}
\end{columns}
\end{frame}

\begin{frame}{Recommended Reading --- Going Further}
\begin{columns}[T]
\column{0.50\textwidth}
\begin{cbox}[ISSAnavy]{Core texts}\scriptsize\begin{tl}\item S. Russell \& P. Norvig, \textit{Artificial Intelligence: A Modern Approach} (4th ed., 2020) --- the AI foundations of L1\item T. Hastie, R. Tibshirani \& J. Friedman, \textit{The Elements of Statistical Learning} (2nd ed., 2009) --- L2--L4 in depth\item C. Bishop, \textit{Pattern Recognition and Machine Learning} (2006) --- probabilistic view; GMM and EM\item I. Goodfellow, Y. Bengio \& A. Courville, \textit{Deep Learning} (2016) --- neural networks and autoencoders\item R. Sutton \& A. Barto, \textit{Reinforcement Learning: An Introduction} (2nd ed., 2018) --- L6--L8\end{tl}\end{cbox}
\column{0.48\textwidth}
\begin{cbox}[AccentCol]{Data interpretation and defence context}\scriptsize\begin{tl}\item D. Huff, \textit{How to Lie with Statistics} (1954) --- the classic on misleading data\item D. Kahneman, \textit{Thinking, Fast and Slow} (2011) --- regression to the mean, base rates, judgement biases\item J. Pearl \& D. Mackenzie, \textit{The Book of Why} (2018) --- correlation, causation and confounding\item P. Scharre, \textit{Army of None} (2018) --- autonomous weapons and human control\end{tl}\end{cbox}
\begin{ibox}\scriptsize The companion reference text for each lecture expands every slide of this course in full detail.\end{ibox}
\end{columns}
\end{frame}
\end{document}
